\chapter{Ricci curvature bounded below}
\source{cheeger-gromov-taylor:page-0010,cheeger-gromov-taylor:page-0011,cheeger-gromov-taylor:page-0012,cheeger-gromov-taylor:page-0013,cheeger-gromov-taylor:page-0014,cheeger-gromov-taylor:page-0015,cheeger-gromov-taylor:page-0016,cheeger-gromov-taylor:page-0017,cheeger-gromov-taylor:page-0018,cheeger-gromov-taylor:page-0019}

Theorem~\ref{thm:kernel-estimates-bounded-curvature} extends to manifolds with boundary after the corresponding
generalization of Proposition~\ref{prop:refined-injectivity-angle-lower-estimates}.

Let $M^n$ be an arbitrary complete Riemannian manifold with $\partial M^n = \emptyset$.
Let $\Tscr_2^{\phi,A}$ be the class of even functions $\{f(\lambda)\}$, such that for some constants
$A$, $k_0$, $s_0$ we have

\begin{equation}\tag{2.1}
\bigl|f^{(k)}(s)\bigr| \le c_1\left(\frac{k}{A}\right)^k \phi(s),\qquad k \ge k_0,\ s \ge s_0,
\end{equation}

where $\phi \ge 0$ belongs to $L^1(a,\infty)$ for all $a > 0$. If we set

\begin{equation}\tag{2.2}
\psi(r) = \int_r^\infty \phi(s)\,ds,
\end{equation}

then

\begin{equation}\tag{2.3}
\int_r^\infty \bigl|f^{(k)}(s)\bigr|\,ds \le c_1\left(\frac{k}{A}\right)^k \psi(r).
\end{equation}

For $\rho(x_1,x_2)=d=2A+r$, Corollary~\ref{cor:derivative-tail-estimate} gives

\begin{equation}\tag{2.4}
\bigl\|\Delta^k f(\sqrt{-\Delta})\Delta^l u\bigr\|_{B_A(x_1)}
\le c\left(\frac{2k + 2l}{A}\right)^{2k+2l}\psi(r)\|u\|_{B_A(x_2)},
\end{equation}

for any $u \in L^2$ with $\supp u \subset B_r(x_2)$. Thus for such $u$, for $|y| \le A/2$, the
power series

\begin{equation}\tag{2.5}
\begin{aligned}
v(x,y) &= \sum_{k=0}^\infty \frac{y^{2k}}{(2k)!}(-\Delta)^k f(\sqrt{-\Delta})u,\\
&= (\cosh \sqrt{-\Delta}\, y)\,f(\sqrt{-\Delta})u(x),\ (\mathrm{formally})
\end{aligned}
\end{equation}

converges to an element of $L^2([ -A/2,\, A/2]\times B_r(x_1))$. Moreover,

\begin{equation}\tag{2.6}
\left(\Delta_M + \frac{\partial^2}{\partial y^2}\right)v = 0,
\end{equation}

and by (2.4),

\begin{equation}\tag{2.7}
\|v\|_{[-A/2,\,A/2]\times B_r(x_1)} \le c\psi(r)\|u\|_{B_r(x_2)}A^{1/2}.
\end{equation}

The harmonic-function estimate used for point evaluation yields

\begin{equation}\tag{2.8}
\left|k_f(u)\right|_{B_{A/2}(x_1)}
\le k_nA^{-n/2}\bigl[\tilde{\omega}(B_A(x_1))\bigr]^{-(n+1)/2}\psi(r)\|u\|_{B_A(x_1)}.
\end{equation}

This implies that for all $x_1$,

\begin{equation}\tag{2.9}
\|k_f(x_1,x_2)\|_{B_A(x_2)}
\le k_nA^{-n/2}\bigl[\tilde{\omega}(B_A(x_1))\bigr]^{-(n+1)/2}\psi(r),
\end{equation}

where $k_n$ is a constant depending only on $n$. Similarly, we obtain

\begin{equation}\tag{2.10}
\|\Delta_2 k_f(x_1,x_2)\|_{B_A(x_2)}
\le k_nA^{-n/2}\bigl[\tilde{\omega}(B_A(x_1))\bigr]^{-(n+1)/2}
\left(\frac{2l}{2A}\right)^{2l}\psi(r).
\end{equation}

Then for each fixed $x_1$ we can set

\begin{equation}\tag{2.11}
w(x_1,x_2,y)=\sum_{l=0}^{\infty}\frac{y^{2l}}{(2l)!}\Delta_2^l f(\sqrt{-\Delta})\delta_{x_2},
\end{equation}

to get a harmonic function on $[-A/2,A/2]\times B_A(x_2)$ whose $L^2$ norm satisfies

\begin{equation}\tag{2.12}
\|w(x_1,x_2,y)\|_{[-A/2,A/2]\times B_A(x_2)}
\le c\|k_f(x_1,x_2)\|_{B_A(x_2)}.
\end{equation}

Applying the same estimate in the second variable gives

\begin{equation}\tag{2.13}
\left|k_f(x_1,x_2)\right|_{B_{A/2}(x_1)\times B_{A/2}(x_2)}
\le c_n^2A^{-n}\bigl[\tilde{\omega}(B_A(x_1))\tilde{\omega}(B_A(x_2))\bigr]^{-(n+1)/2}\psi^2(r)
\end{equation}

for $\rho(x_1,x_2)=d=2A+r$. The same argument shows that for all $x_1,x_2$,

\begin{equation}\tag{2.14}
\left|k_f(x_1,x_2)\right|_{B_{A/2}(x_1)\times B_{A/2}(x_2)}
\le c_n^2A^{-n}\bigl[\tilde{\omega}(B_A(x_1))\tilde{\omega}(B_A(x_2))\bigr]^{-(n+1)/2}\psi^2(0),
\end{equation}

provided $\varphi\in L^1(\R^+)$. Note that if $k_{f_\lambda}$ denotes the kernel of
$f(\sqrt{-\Delta})$ with respect to the metric $g$, then
$k^g_{f(\lambda)}=c^{-n/2}_g k_{f(\lambda/c)}$. Thus using (2.13), (2.14) we
can estimate $k^g_{f(\lambda)}$ in two different ways. It is straightforward to
check that the estimates so obtained coincide; i.e., (2.13), (2.14) have the
correct scaling property.

Proposition~\ref{prop:refined-injectivity-angle-lower-estimates} gives a lower estimate for
$\tilde{\omega}(A,x_0)$, if $A$ is sufficiently small. To simplify the statement
of Theorem~\ref{thm:kernel-estimates-ricci-lower-bound}, we will use only Proposition~\ref{prop:refined-injectivity-angle-lower-estimates}(i) even though (ii) is sharper
if $\rho(p,x_j)=\rho_j$ is sufficiently large. Define $r(H,V)$ by

\begin{equation}\tag{2.15}
V_r^H(H,V)=\frac{V}{2},
\end{equation}

where $V_r^H$ is given by (1.27).

\begin{theorem}[Kernel estimates with Ricci lower bound]
\label{thm:kernel-estimates-ricci-lower-bound}
\source{cheeger-gromov-taylor:page-0012}
\dcref{Theorem 2.1}
\group{section-2}

\textit{Let $M^n$ be a complete Riemannian manifold with $\partial M^n=\emptyset$, and $\operatorname{Ric}_M \ge (n-1)H$. Set $V_{r_0}=V_{r_0}(p)$. If $f \in \mathcal{F}^{\varphi,A}_2$ and $A < \frac12 r(H,V)$, then the kernel $k_f(x_1,x_2)$ is continuous for $x_1 \neq x_2$ and satisfies (2.13), (2.14), where}
\begin{equation}\tag{2.16}
\widetilde{\omega}(B_A(x_j)) \ge \frac{1}{\left(2V_{r+r_0}^{H}/V_{r_0}\right)-1}.
\end{equation}
\end{theorem}

\begin{remark}
\label{rem:gradient-kernel-estimate}
\source{cheeger-gromov-taylor:page-0012}
\dcref{Remark after Theorem 2.1}
\group{section-2}

Under the hypotheses of Theorem~\ref{thm:kernel-estimates-ricci-lower-bound}, domination of semigroups also gives
pointwise bounds for the norm of the gradient of $k_f$.
\end{remark}

The heat kernel is obtained from Theorem~\ref{thm:kernel-estimates-ricci-lower-bound}
by taking $f_t(\lambda)=e^{-t\lambda^2}$; its Gaussian decay is controlled by
the same volume and injectivity-angle bounds.

For dimensions $2$ and $3$, the same pointwise estimates follow assuming
$\Ric_M\geq(n-1)H$ and $V_{r_0}(p)\geq V>0$.

Let
\begin{equation}\tag{2.35}
\Hg=e_i e_j(g)-\nabla_{e_i}\nabla_{e_j}(g)
\end{equation}
denote the Hessian of $g$. The coercive estimate is
\begin{equation}\tag{2.36}
\|Hg\|_{B_{r/2}(x_0)}^{2}\le c(r)\Bigl(\|dg\|_{B_r(x_0)}^{2}+\|\Delta g\|_{B_r(x_0)}^{2}\Bigr)
\end{equation}
when $\Ric_M\ge H$. A lower bound for the refined injectivity angle then
controls the Sobolev constant; combined with (2.36), this gives the pointwise
kernel estimate in dimensions $2$ and $3$ for $f\in\mathscr{F}_1^1$.

For a $1$-form $\omega$, use the Bochner-Weitzenbock formula
\begin{equation}\tag{2.37}
(d\delta+\delta d)\omega=\sum_i\bigl(-\nabla_{e_i}\nabla_{e_i}+\nabla_{\nabla_{e_i}e_i}\bigr)\omega+\Ric(\omega^*),
\end{equation}
where
\begin{equation}\tag{2.38}
\Ric(\omega^*)=(\Ric(\omega))^*,
\end{equation}
and $\omega^*$ is the tangent vector dual to $\omega$. Fix a metric ball $B_r(x)$, and let $\psi(r)$ be defined by $\psi(r)|_{[0,r/2]}\equiv 1$, $\psi(r)|_{[r/2,r]}\equiv -2r/\tilde r+2$. By Stokes' theorem,
\begin{equation}\tag{2.39}
\int_{B_{\tilde r}(x_0)}\psi^2\, d\delta dg\wedge *dg
=\int_{B_{\tilde r}(x_0)}\psi^2\delta dg\wedge *\delta dg
-\int_{B_{\tilde r}(x_0)}2\psi\, d\psi\delta dg\wedge *dg.
\end{equation}

Applying (2.37) to the left-hand side and integrating by parts gives
\[
\int_{B_r(x_0)} \psi^2 \delta \dg \wedge *\dg
= \int_{B_r(x_0)} \left\{ \psi^2 \sum_i \left(-\nabla_{e_i}\nabla_{e_i} + \nabla_{\nabla_{e_i} e_i}\right)\dg \wedge *\dg + \psi^2 \Ric(\dg) \wedge *\dg \right\}
\]
\begin{equation}
= \int_{B_r(x_0)} \left\{ \psi^2 \sum_i \nabla_{e_i}\dg \wedge *\nabla_{e_i}\dg + 2\psi \sum_i e_i(\psi)\nabla_{e_i}\dg \wedge *\dg
+ \psi^2 \Ric(\dg) \wedge *\dg \right\}.
\tag{2.40}
\end{equation}

If we substitute (2.40) into (2.39) and transpose the middle term we get
\[
\int_{B_r(x_0)} \left( \psi^2 \sum_i \nabla_{e_i}\dg \wedge *\nabla_{e_i}\dg + \Ric(\dg)\wedge *\dg \right)
\]
\begin{equation}
= \int_{B_r(x_0)} \psi^2 \delta \dg \wedge *\delta \dg
- \int_{B_r(x_0)} 2\psi\, d\psi\,\delta \dg \wedge *\dg
- \int_{B_r(x_0)} 2\psi \sum_i e_i(\psi)\nabla_{e_i}\dg \wedge *\dg.
\tag{2.41}
\end{equation}

Note that $\nabla \dg = Hg$. Applying the Schwarz inequality to the last two terms of the right-hand side of (2.41) gives
\[
\int_{B_r} \psi^2 \sum_i \nabla_{e_i}\dg \wedge *\nabla_{e_i}\dg + (n-1)H \|\psi\,dg\|^2_{B_r(x_0)}
\]
\begin{equation}
\begin{aligned}
&\le \|\Delta g\|^2_{B_r(x_0)}
+ \left(\|\Delta g\|_{B_r(x_0)} + \frac{4}{r^2}\|dg\|^2_{B_r(x_0)}\right) \\
&\quad + \left(\frac12\int_{B_r(x_0)} \psi^2 \sum_i \nabla_{e_i}\dg \wedge *\nabla_{e_i}\dg
+ \frac{8}{r^2}\|dg\|^2_{B_r(x_0)}\right).
\end{aligned}
\tag{2.42}
\end{equation}

If we transpose the fourth term on the right-hand side, we obtain

\begin{proposition}[Coercive Hessian estimate]
\label{prop:coercive-hessian-estimate}
\source{cheeger-gromov-taylor:page-0015,cheeger-gromov-taylor:page-0016}
\dcref{Proposition 2.2}
\group{section-2}

\textit{Let $M^n$ be complete, and $\Ric_M \geq (m-1)H$. Then for all $x, r$}
\begin{equation}\tag{2.43}
\|Hg\|^2_{B_{r/2}(x)} \le 4\|\Delta g\|^2_{B_r(x)} + \frac{24}{r^2}\|dg\|^2_{B_r(x)}
- 2(n-1)H\|\psi\,dg\|^2_{B_r(x)}.
\end{equation}
\end{proposition}

Let $A(r,\theta), \mathcal{Q}^H(r)$ be as in (1.28). Given $\epsilon > 0$ and $H$, we define a refinement of the injectivity angle by letting $\widetilde{\omega}(H,\epsilon,r,x)$ be the angular
measure of the set $\Omega^H(\epsilon)$ of initial tangent vectors to geodesics $\gamma$ from $x$ such
that $\gamma|_{[0,r]}$ is minimal and

\begin{equation}\tag{2.44}
\frac{A(s,\theta)}{\mathcal{Q}^H(s)} \ge \epsilon,\qquad 0 \le s \le r,
\end{equation}

at $\gamma(s)$. Note that for all $H$, $\widetilde{\omega}(H,0,r,x)=\widetilde{\omega}(r,x)$. Let $\nabla^i g$ denote the $i$-th
covariant derivative of $g$.

\begin{proposition}[Pointwise Sobolev estimate]
\label{prop:pointwise-sobolev-estimate}
\source{cheeger-gromov-taylor:page-0017}
\dcref{Proposition 2.3}
\group{section-2}

\textit{Let $g$ be a smooth function on $B_r(x_0)\subset M^n$. Then for}
$N>n/2$,
\begin{equation}\tag{2.45}
g^2(x_0)\le \frac{c_3(N)}{\epsilon\,\widetilde{\omega}(H,\epsilon,\bar r,x_0)}
\left(\int_0^{\bar r} r^{2(N-1)}\,\frac{dr}{\mathcal{Q}^H(r)}\right)
\sum_{i=0}^N \bar r^{2(i-N)}\|\nabla^i g\|^2_{B_{\bar r}(x_0)}.
\end{equation}
\end{proposition}

\begin{proof}
\source{cheeger-gromov-taylor:page-0017}
\sketch
Let $\psi$ be essentially as above, but made $C^\infty$. Then
\begin{equation}
g^2(x_0)\le \frac{1}{\widetilde{\omega}(H,\epsilon,\bar r,x_0)}
\int_{\Omega^H(\epsilon)}\left|\int_0^{\bar r}\frac{d}{dr}\bigl(\psi g(r,\theta)\bigr)\,dr\right|^2 d\theta
\end{equation}
\begin{equation}\tag{2.46}
= \frac{1}{\widetilde{\omega}(H,\epsilon,\bar r,x_0)}
\cdot \int_{\Omega^H(\epsilon)}\left[\int_0^{\bar r}\frac{r^{N-1}}{A^{1/2}(r,\theta)}\frac{d^N}{dr^N}\bigl(\psi g(r,\theta)\bigr)A^{1/2}(r,\theta)\,dr\right]^2 d\theta
\end{equation}
\[
\le \frac{c_3(N)}{\epsilon\,\widetilde{\omega}(H,\epsilon,\bar r,x_0)}
\left(\int_0^{\bar r}\frac{r^{2(N-1)}\,dr}{\mathcal{Q}^H(r)}\right)
\sum_{i=0}^N \bar r^{2(i-N)}\int_{B_{\bar r}(x_0)}\|\nabla^i g\|^2,
\]
and the claim follows.
\end{proof}

Proposition~\ref{prop:refined-injectivity-angle-lower-estimates} gives a lower estimate for $\widetilde{\omega}(H,\epsilon,\bar r,x_0)$ for $\epsilon,\bar r$
sufficiently small. Combining this with Proposition~\ref{prop:coercive-hessian-estimate} and
Proposition~\ref{prop:pointwise-sobolev-estimate} gives
get the following result. As in previous instances, the statement can be
sharpened somewhat; see Proposition~\ref{prop:refined-injectivity-angle-lower-estimates}(iv).

\begin{theorem}[Ricci-lower-bound pointwise estimate]
\label{thm:ricci-lower-bound-pointwise-estimate}
\source{cheeger-gromov-taylor:page-0017,cheeger-gromov-taylor:page-0018}
\dcref{Theorem 2.4}
\group{section-2}

\textit{Let $M^n$ be complete, and $\mathrm{Ric}_M\ge (n-1)H$. Fix $p, x\in M^n$}
\textit{and $r_0>0$. Set $V_{r_0}=V_{r_0}(p)$, $\rho=\rho(p,x)$, $V^H_{r(H,V_{r_0})}=\frac12 V_{r_0}$.}

\noindent (i) Then for $N>n/2$,
\begin{equation}\tag{2.47}
g^2(x)\le \frac{c_3(N)}{\bigl[(4V^H_{\rho+r_0}/V_{r_0})-3\bigr]^2}
\left(\int_0^{r(H,V_{r_0})}\frac{r^{2(N-1)}\,dr}{\mathcal{Q}^H(r)}\right)
\sum_{i=0}^N \bar r^{2(i-N)}\|\nabla^i g\|^2_{B_r(x)}.
\end{equation}

\noindent (ii) In particular, for $n=2,3$,
\begin{equation}\tag{2.48}
\begin{aligned}
g^2(x)&\le \frac{c_3(2)}{\bigl[(4V^H_{\rho}/V_{r_0})-3\bigr]^2}
\left(\int_0^{r(H,V_{r_0})}\frac{r^2}{\mathcal{Q}^H(r)}\,dr\right) \\
&\quad\times\Bigl[r_0^{-4}\|g\|^2_{B_{r_0}(x)}
+25r_0^{-2}\|dg\|^2_{B_{r_0}(x_0)} \\
&\qquad -2(n-1)H\|\psi dg\|^2_{B_{r_0}(x_0)}
+4\|\Delta g\|^2_{B_{r_0}(x_0)}\Bigr].
\end{aligned}
\end{equation}
\end{theorem}

In view of (2.48), the kernels $k_f$ for $f\in \mathscr{F}_1^1$ satisfy the same
type of pointwise control.
